\section{Financial validation and interpretation}

A numerical stopping problem can be well defined while its market interpretation is unresolved. Empirical application requires dated contract terms, a justified ownership state, cash information available at the valuation cutoff, and funding appropriate to the payment obligations. Event and source timestamps may describe different objects. Simultaneous records can require an explicit ordering rule; absent or stale inputs must not silently become valid observations.

Joint call and put compatibility under a shared state, cash specification, funding convention and volatility is an informative diagnostic. Separate unconstrained calibrations can remove that restriction. Compatibility with bid--ask intervals is still not identification of the correct model or an executable arbitrage. A model that fails a renewed quote diagnostic after a sensitivity change can remain a conditional mathematical case, but must not be reclassified as an accepted empirical specification because its policy loss is small.

A quoted spread offers a monetary scale. It is not automatically the cost of exercise, hedging or implementation. An upper loss limit below the spread establishes only that conditional scale comparison. A limit above it leaves the inequality unestablished; it does not prove that true loss exceeds the spread.

A prospective evaluation should freeze its observations, admissibility rule, calibration and inference before outcomes can affect method choice. If no observation passes the prescribed screen, a performance estimate is unavailable. Wider post-sample thresholds can inform a future design but cannot silently replace the primary sample. Numerical replication, admissible measurement and independent economic validation are different requirements.

\section{Conclusions and data availability}

I distinguish policy quality from the tightness of its evaluation bound and from the financial opportunities used as the benchmark. Analytically centred values can improve the measurement of loss without changing the policy. A larger calendar changes the comparison problem, while a feasible terminal adjustment changes the rule itself. Common-upper comparisons with valid martingales identify expected policy gains; finite-family bounds and explicit timing accounting delimit what the resulting evidence can support.

I use this separation to distinguish two allocation questions. Additional computation spent on the dual construction can improve knowledge of a fixed policy's loss; additional computation spent on training can change the policy's reward. Neither expenditure can be ranked solely by the resulting dual gap because that gap combines loss and slack. A comparison of training and evaluation budgets therefore requires a common financial problem, a declared computational workload and independent evaluation of the resulting policies. The framework specifies that comparison but does not identify an optimal allocation without evidence on both margins.

Predictable compensation identifies a separate source of error if the conditional mean-zero requirement is relaxed. A zero apparent gap or a common upper payoff cannot certify policy quality in its presence. I therefore retain analytical centring as a validity condition, separate from the approximation choices used to tighten the bound.

These distinctions provide a basis for a conditional computational investigation. They do not establish which specification performs best in observed markets or how the results generalise to additional state variables and uncertain funding or cash. Such conclusions require their own numerical and empirical evidence.

In this methods-only publication, I set out the model and apply the cited methods. The finite-tree example is stipulated for analytical verification. I include no raw or cleaned subscription observations, fitted market inputs, empirical summaries, research-result tables or data-derived figures. I retain the complete numerical paper separately. The LaTeX sources and build instructions for this version are available in the accompanying repository.
